\begin{abstract}
Indonesian employers withhold employee income tax (PPh 21) monthly under rules that changed 13 times in 2023--2026, five of them retroactively. Calculators that hard-code these rules need new code for each change and cannot reproduce earlier results after a retroactive rule appears. This study presents an adaptive knowledge-based system that keeps these rules current with a large language model (LLM). Government regulations and each company's pay policies form two layers of one knowledge graph in which every rule records when it is valid and when it became known. A deterministic engine computes the tax from these records and recalculates paid months when a retroactive rule arrives. The LLM only drafts updates, each of which must quote its source and pass deterministic gates. The system abstains when the drafts disagree. The evaluation replayed the documents of five real companies with 13 regulations, 6 labor agreements, and 20 distractors and scored 26,328 payslips against a reference checked by a tax practitioner. With a human-capital software engineer approving the proposals, payslip accuracy on the payroll date was 0.945 to 1.000 in four companies and exceeded that of a code-based update made 26 days after each issue. An incomplete labor-agreement installation lowered the fifth company to 0.239, so the pre-specified pooled hypothesis failed. A completeness gate added after the evaluation removed this failure in a replay of the recorded drafts. The engineer wrote updates by hand for 11 of 138 documents. Every late upload of a retroactive rule left the earlier payroll reproducible.
\end{abstract}

\begin{keywords}
Bitemporal knowledge graph, knowledge acquisition, knowledge-based systems, large language models, payroll tax, retroactive regulation.
\end{keywords}
